\chapter{Related work} \label{sec:related_work}

As sparse kernels are becoming pervasive, there is a large body of work around performance analysis and optimizations---both hardware and software---of specific sparse kernels such as SpMV, and others. However, none of these approaches provides a comprehensive and flexible solution that works efficiently for general sparse tensor algebra. This chapter reviews such related work, structured around the core contributions of this thesis.

\section{Performance Analysis of Sparse Workloads}

Because of its unique characteristics, sparse computation has been widely studied in the past. However, existing performance analyses mostly focus on SpMV kernels, and very little work has been dedicated to other scientific kernels or other applicaitons like sparse machine learning.

\subsection{Performance Analysis of SpMV}
SpMV is one of the most critical kernels in scientific computing, and its performance has been extensively studied, particularly on CPUs. Mantovani et al.~\cite{mantovani2020benchmark} evaluate thread- and node-level scalability on the MareNostrum4 supercomputer and ThunderX2-based systems, comparing both performance and energy consumption across microarchitectural choices such as vector unit width and cache hierarchy design. Armejach et al.~\cite{armejach2019hpc} analyze a range of HPC processors employing different combinations of in-order and out-of-order cores as well as varying last-level cache capacities per core, concluding that scaling last-level caches yields diminishing returns while substantially increasing area and power. Oliveira et al.~\cite{oliveira2021damov} compare HPC processors using metrics including GFLOPS per Watt, GFLOPS versus memory bandwidth, and bytes per FLOP. Other studies have focused on specific characterization metrics: Bean et al.~\cite{bean2015gdma} characterize data movement in memory-intensive benchmarks (e.g., MPKI, last-to-first-level cache ratios, operational intensity) and examine the microarchitectural trade-offs underlying these behaviors. Sparse matrix formats have also received attention. Asgari et al.~\cite{asgari2021copernicus} discuss the performance implications of different sparse matrix representations, while Alappat et al.~\cite{alappat2022spmv} evaluate SpMV and lattice quantum chromodynamics kernels on the A64FX platform, comparing the impact of multiple formats. As GPUs have become increasingly prominent in HPC, several authors have investigated SpMV performance on them~\cite{bell2009spmv}. More recently, Nguyen et al.~\cite{fpga2022ccpe} analyze SpMV performance on off-the-shelf CPUs, GPUs, and FPGAs, demonstrating that indirect lookups into the dense vector remain a dominant bottleneck across all three architectures.

\subsection{Performance Analysis of Other Sparse Kernels}
Despite the extensive body of work on SpMV, comparatively little effort has been devoted to examining more complex sparse kernels such as SpADD or SpGEMM~\cite{gao2023spgemm_survey}, which introduce additional challenges---including more irregular control flow, more complex merging behavior, and heavier reliance on intermediate data structures.

\subsection{Performance Analysis of Embedding Operations}
With the rise of sparsity and feature embeddings in machine learning, several authors have begun examining their architectural implications. Gupta et al.~\cite{dlrm2020hpca} show that the embedding bag function in DLRMs suffers from memory bottlenecks analogous to scientific kernels, despite it has different spatio-temporal locality characteristics. Zaheer et al.~\cite{bigbird2024neurips} highlight the inefficiency of loading sparse attention blocks in LLMs on conventional architectures. However, none of these studies analyze performance at the microarchitectural level, nor do they characterize these workloads in ways that provide actionable guidance for designing more efficient architectures.

\subsection{Our Solution}
To overcome these limitations, in \Cref{sec:architectural-implications-of-sparsity} we provide a comprehensive characterization of sparsity in both scientific and machine learning models. We start by providing, in \Cref{sec:architectural-implications-of-sparsity-scientific-computing}, a comprehensive performance analysis of a broad suite of sparse tensor algebra kernels in scientific computing, including SpMV, SpGEMM, and SpADD, evaluated on a large set of matrices and diverse processor architectures. Then, in \Cref{sec:architectural-implications-of-sparsity-in-machine-learning} we extend this characterization to a large set of embedding operations, and use the results of this analysis to define the ideal architectural properties for these workloads.

\section{Architectures Optimized for Sparse Computation}

Because of their importance, several authors proposed different architectures to accelerate sparse workloads. While most authors focus on CPU optimizations for indirect memory accesses, some authors also propose custom architectures for specific sparse kernels and general sparse tensor algebra kernels. However, conversely from our DAE architecture, these solutions offer either limited performance improvement or flexibility.

\subsection{Architectural Optimizations for Sparse Tensor Algebra}
A number of solutions have been proposed to optimize indirect lookups on traditional architectures. Programmable prefetchers~\cite{yu2015, AinsworthJ16, talati2021prodigy} allow software to describe a restricted set of access patterns for the hardware to prefetch, including the indirect access in SpMV. Data layout transformation mechanisms, instead, reorganize sparse data into denser formats to improve locality and better exploit the memory hierarchy~\cite{dre, planar2021acm, impulse}. Vector extensions like Tarantula~\cite{tarantula} extend support for non-unit-stride accesses and improves the efficiency of gather/scatter instructions. One of the main limitations of these proposals is to only support limited classes of access patterns---largely those arising from SpMV---and cannot accelerate general sparse tensor algebra. Moreover, they do not address data-dependent control-flow overheads. A more general solution is the streaming abstraction proposed by Scheffler et al.~\cite{scheffler2021stream,scheffler2023stream,scheffler2024stream,scheffler2025occamy,zhang2023axipack,manoni2025spikestream}. The authors extend the RISC-V instruction set to support streaming abstractions for indirect accesses and merging. The core can program these streams to load and merge sparse/dense tensors, and just pop the resulting data to be computed. This offloads indirect accesses and control flow from the compute pipeline, achieving substantial performance improvements. While this solution can accelerate a larger set of sparse algebraic kernels, similarly to the other optimizations in the literature, it mostly targets scalar cores, not state-of-the-art vector cores, leaving performance on the table.

\subsection{Architectures for Sparse Tensor Algebra}
At the opposite end of the design spectrum lie fully custom architectures for sparse tensor algebra. Several architectures have been proposed to accelerate SpMV-like workloads, both on FPGAs~\cite{fowers2014spmv,grigoras2015spmv,umuroglu2014spmv}, fixed-function architectures~\cite{nurvitadhi2015spmv,zhang2020sparch,sadi2019spmv}, and processing in memory systems~\cite{giannoula2022sparsep}. Designs such as GAMMA~\cite{zhang2021gamma} and related systems~\cite{srivastava2020matraptor,pal2018outerspace,zhang2020sparch,lin2010spmv,quin2020sigma,parashar2017sparse,zhang2016cambriconx} implement fixed-function units for sparse matrix--matrix multiplication, a fundamental kernel to which many tensor algebra operations can be reduced. These solutions offer high performance but very limited flexibility and no opportunities for kernel fusion~\cite{bharadwaj2022sparse_for_ml}.

\subsection{Architectures for General Sparse Tensor Algebra}
More flexible architectures, such as Capstan~\cite{rucker2021capstan} and others~\cite{dadu2019spu,hegde2019extensor,asgari2020alrescha,hwang2020centaur,mishra2017sparse,nurvitadhi2016sparse}, can be programmed to execute a wider range of Einsum expressions~\cite{sam2023asplos}. Despite their broader applicability, we argue that (i) computation on specialized hardware is constrained by design choices such as a restricted set of operations, data types, and memory capacity, and (ii) using such accelerators typically requires moving tensors into and out of accelerator memory before and after every computation. In real-world applications such as CP-ALS tensor decomposition~\cite{phipps2019genten}, where intermediate results must be evaluated at each iteration, this overhead can outweigh the benefits.

\subsection{Hardware-Software Co-Design and DAE Architectures}
To address these limitations, DAE architectures combine specialized access units with general-purpose compute cores. Maple~\cite{maple2022isca} uses specialized access units to accelerate indirect memory accesses for a scalar core. Similarly, HATS~\cite{hats} and SpZip~\cite{spzip2021isca} employ dedicated access units to accelerate graph traversals: HATS exploits graph locality to reduce data movement, while SpZip applies compression and decompression techniques to lower memory traffic. Instead, Smash uses a custom engine to decompress sparse matrices stored in a custom compressed format~\cite{kanellopoulos2019smash}. Cambricon-S uses a specialized engine to efficiently load weights of sparse neural networks~\cite{zhou2018cambricons}. Unlike programmable prefetchers, which aim to bring data into caches ahead of time, these access units push data directly into hardware or software queues that the compute core consumes, making them more effective for highly irregular access patterns. However, these systems are not designed to support general sparse tensor algebra, and are restricted to scalar cores rather than modern vectorized architectures.

\begin{sidewaystable}
    %\scriptsize
    \centering
    \caption{Qualitative state-of-the-art comparison.}
    \vspace{-0.5em}
    \label{tab:soa-comparison}
    %\vspace{-0.3cm}
    %\resizebox{\textwidth}{!}{%
    \begin{tabular}%{@{}l@{\hspace{-4pt}}c@{\hspace{4pt}}c@{\hspace{4pt}}c@{\hspace{4pt}}c@{\hspace{4pt}}c@{}}
    {@{}ccccccc@{}}
    \toprule
    \textbf{Features} & 
    \textbf{\specialcell{Programmable \\ prefetchers \\ \cite{yu2015, AinsworthJ16, talati2021prodigy}}} &
    \textbf{\specialcell{Gather-scatter \\ \cite{spert2-T0,tarantula}}} &
    \textbf{\specialcell{Traversal \\ engines \\ \cite{hats,spzip2021isca}}}  &
    \textbf{\specialcell{Fixed-function \\ DSAs \\ \cite{zhang2021gamma,srivastava2020matraptor,pal2018outerspace,zhang2020sparch}}} &
    \textbf{\specialcell{Programmable \\ DSAs \\ \cite{rucker2021capstan,dadu2019spu,hegde2019extensor}}} &
    \textbf{TMU}\\ \toprule

\specialcell{Improves memory \\ bandwidth utilization}                  & \vm & \vm & \vm & \vm & \vm & \vm \\
\midrule
\specialcell{Hardware-accelerated \\ tensor traversal}                  &  ~  &  ~  & \vm & \vm & \vm & \vm \\
\midrule
\specialcell{Hardware-accelerated \\ tensor merging}                    &  ~  &  ~  &  ~  & \vm & \vm & \vm \\
\midrule
\specialcell{Sparse and dense tensor \\ algebra complete}               & \vm & \vm &  ~  &  ~  & \vm & \vm \\
\midrule
\specialcell{Compressed storage \\ format complete}                     & \vm & \vm &  ~  &  ~  & \vm & \vm \\
\midrule
\specialcell{General-purpose compute \\ capabilities}                   & \vm & \vm & \vm &  ~  &  ~  & \vm \\
\midrule
\specialcell{SIMD-friendly compute \\ (e.g. compatible with Arm SVE)}   &  ~  & \vm &  ~  & \vm & \vm & \vm \\

\bottomrule
\end{tabular}
%}
%\vspace{-0.4cm}

\end{sidewaystable}



\subsection{DAE Architectures in Production}
Similar DAE designs that separate memory access from computation are becoming increasingly prevalent among hyperscalers. Meta MTIAs~\cite{mtia2023isca} and Google TPUs~\cite{tpu2023isca} both adopt such an approach to accelerate embedding lookups in recommendation models. Recent Nvidia GPUs~\cite{nvidiah100} similarly include specialized access units, Tensor Memory Accelerators (TMAs), for fetching dense tensor tiles. While these designs demonstrate the value of separating access from execution, to the best of our knowledge, none of these architectures can accelerate general sparse tensor algebra. 

\subsection{Our Solution}
We designed the TMU to close these gaps. As discussed in \Cref{sec:the-tmu}, conversely from other approaches, a multicore DAE processor composed TMU-CPU DAE cores achieves an unprecedented combination of efficiency and generality. \Cref{tab:soa-comparison} qualitatively compares the TMU with state-of-the-art solutions.

\section{DAE Compilation}

One of the primary barriers to widespread adoption of DAE architectures is programmability. Scientific applications often rely on kernel libraries (e.g., BLAS), making hand-written and hand-optimized kernels feasible. However, this approach is less suitable for machine learning.

\subsection{Traditional Machine Learning Compilers}
Traditional machine learning compilers generate machine code from a semantic representation of the model~\cite{tf2015whitepaper,pytorch2019neurips,caffe2014arxiv,jax2018github,pytorchgeometric2019github,glow2018arxiv}. As in the Meta MTIA ecosystem~\cite{mtia2023isca}, models are decomposed and lowered to handwritten kernels optimized for the target architecture. While effective, this strategy does not scale to the wide variety of embedding operations, their algorithmic complexity, or the combinatorial explosion of input formats, reduction operators, and algorithmic optimizations (e.g., operator fusion).

\subsection{Tensor Compilers}
Tensor compilers address this challenge by progressively lowering machine learning models through a sequence of IRs until code generation. For example, through MLIR~\cite{mlir2021cgo}, PyTorch programs can be lowered via IRs such as Linalg (i.e., linear algebra kernels like matrix multiplication), SCF (i.e., structured imperative code composed of loops, if-then-else constructs, and memory/compute instructions), and LLVM (i.e., unstructured low-level imperative code composed of branching instructions and memory/compute operations). This enables optimizations at multiple abstraction levels. High-level IRs (e.g., Linalg~\cite{mlir2021cgo}, index notations~\cite{taco2017oopsla,mlirsparse2022taco}, or Einsums~\cite{tf2015whitepaper}) describe what to compute, whereas lower-level IRs describe how to compute it. Although high-level tensor IRs support powerful algorithmic optimizations such as operation fusion~\cite{tf2015whitepaper,pytorch2019neurips,caffe2014arxiv,jax2018github,pytorchgeometric2019github,glow2018arxiv}, their level of abstraction makes them unsuitable for architectural optimizations, including those required for DAE systems. Consequently, most architectural optimizations occur in lower-level IRs (e.g., SCF or LLVM IR~\cite{llvm2004cs}) after high-level algorithmic transformations have already been applied. While tensor compilers have been successfully used for CPUs~\cite{taco2017oopsla,mlirsparse2022taco,looplets2023cgo,mosaic2023pldi,indexstreams2023pldi}, GPUs~\cite{tacoscheduling2020oopsla,sparsetir2023asplos,liu2024unisparse}, FPGAs~\cite{liu2024unisparse}, and distributed systems~\cite{spdistal2022sc}, compilation for tensor DAE systems remains largely unexplored. Lopoukhine et al.~\cite{rvx2025lopoukhine} recently showed the potential of multi-level compilation flows for ISA extensions such as streaming registers, but this technique does not directly apply to complex DAE systems like a TMU--CPU core, where both indexing and memory accesses are fully offloaded to the access unit.

\subsection{DAE Compilation}
Earlier work on DAE compilation~\cite{desc2015micro}, predating modern tensor compilers, explored generating DAE code from LLVM IR~\cite{llvm2004cs}. However, LLVM IR is unstructured and represents loops using conditional branches, making it ill-suited for analyzing the deep loop hierarchies that arise in embedding operations~\cite{panda2020tcad,rvx2025lopoukhine}. For example, identifying loop-offloading candidates (\Cref{sec:lowering-scf-to-slc}) becomes impractical in such representations. High-level synthesis tools address similar challenges using pragmas~\cite{cong2011autoesl,stream2019wang}, but these require user intervention and are incompatible with our automated compilation flow. For affine loops, Jimborean et al.~\cite{dvfs2018jimborean} use polyhedral analysis to overcome these limitations, but this approach cannot be applied to embedding operations. Moreover, unlike our SLC IR, LLVM IR is not designed to represent DAE programs: it cannot express separate access and compute stages, nor can it model data (de)serialization. As a result, optimizations cannot be applied iteratively and must instead be embedded directly into the lowering and code-generation process, which is impractical for sophisticated transformations.

\subsection{Our Solution}
To overcome these limitations, in \Cref{sec:dae-compilation} we introduce the Ember compiler, which features custom IRs to lower PyTorch and TensorFlow embedding operations to optimized DAE code, enabling the full capabilities of DAE architectures at no programmability cost.

\section{Summary}
In the end, we believe that current optimizations and architectures for sparse tensor algebra either focus on performance or generality, but not both. This thesis bridges this gap.